\chapter{Probability in Machine Learning}

There is a temptation, when you first meet machine learning, to picture it as a
bag of clever tricks: a decision tree here, a neural network there, a gradient
descent loop turning the crank. Useful, yes -- but it leaves you with a hundred
disconnected recipes and no story. The goal of this chapter is to give you the
story. And the story is short enough to state in one sentence:

\begin{focusbox}
Almost every supervised and unsupervised learning method you will ever meet is,
underneath, a \emph{probabilistic model} together with a rule for choosing its
parameters from data. Learning \emph{is} estimation. Prediction \emph{is}
computing a conditional distribution.
\end{focusbox}

Once you believe that sentence, machine learning stops being a zoo and becomes a
single landscape with familiar landmarks. Maximum likelihood (Chapter on
estimation), Bayes' theorem (the conditioning chapter), the normal distribution,
the law of large numbers, KL divergence from information theory -- they all
reappear here wearing work clothes. This is the flagship application chapter: it
ties the book together. We will move slowly, build each idea from a concrete
story, and stay honest about what works, what fails, and why.

\section{The likelihood view of learning}
\label{ch34:likelihood-view}

\subsection{A story: the biased coin and the spam filter}

Suppose a friend hands you a bent coin and asks, ``what's the probability it
lands heads?'' You flip it $20$ times and see $14$ heads. You would not hesitate:
your gut says about $14/20 = 0.7$. Where does that number come from? Not from a
formula you memorized -- from a deeper instinct. You are reaching for the value of
$p$ that makes the data you actually saw \emph{least surprising}. A coin with
$p=0.7$ produces ``$14$ out of $20$'' far more readily than a coin with $p=0.1$
or $p=0.95$. You are, without naming it, doing maximum likelihood.

Now zoom out. A spam filter sees an email and must decide \texttt{spam} or
\texttt{ham}. A medical model sees a scan and outputs a probability of disease. A
recommender sees your clicks and predicts the next one. In every case there is an
input $x$, an output $y$, and a family of probability distributions
$p(y \given x; \theta)$ indexed by parameters $\theta$. \emph{Learning} is
choosing $\theta$ from a training set. \emph{Prediction} is plugging in a new $x$
and reading off the distribution over $y$. The bent coin is the whole subject in
miniature.

\begin{intuitionbox}
A model is a \emph{machine for assigning probabilities to data}. A good model
assigns high probability to the things that actually happened and low
probability to things that did not. Training tunes the knobs $\theta$ to do
exactly that on the data you have, in the hope it generalizes to data you do not.
\end{intuitionbox}

\subsection{Datasets, models, and the i.i.d.\ assumption}

Fix the setup we will use all chapter. We observe a dataset
\[
  \mathcal{D} = \set{(x_1,y_1),\dots,(x_n,y_n)},
\]
$n$ examples, each a pair of an input (or \emph{feature vector}) $x_i$ and a
label $y_i$. We posit a parametric model: a family of distributions
$p(\,\cdot \given \theta)$. The single most common assumption -- so common it is
often left unstated -- is that the examples are drawn \emph{independently from one
fixed distribution}, i.i.d.:
\[
  (x_i,y_i) \overset{\text{i.i.d.}}{\sim} p^\star, \qquad i=1,\dots,n.
\]
Here $p^\star$ is the true, unknown data-generating distribution. We never see
$p^\star$ directly; we see $n$ samples from it and try to approximate it with our
model $p(\,\cdot\given\theta)$.

\begin{trapbox}
The i.i.d.\ assumption is a \emph{modeling choice}, not a law of nature. Time
series are not i.i.d.\ (today depends on yesterday). Data with distribution
shift -- the test set drawn from a different $p^\star$ than the training set --
breaks every guarantee in this chapter. When a model fails mysteriously in
production, ``the i.i.d.\ assumption was violated'' is the first suspect.
\end{trapbox}

\subsection{The likelihood and its logarithm}

The \emph{likelihood} of the parameters $\theta$ given the data is the
probability the model assigns to the data, read as a function of $\theta$:
\[
  L(\theta) = p(\mathcal{D}\given\theta) = \prod_{i=1}^n p(x_i,y_i\given\theta),
\]
the product because of independence. We choose $\theta$ to make this large. But a
product of $n$ small numbers underflows to zero on any real computer, and
products are awkward to differentiate. So we take logarithms. Because $\log$ is
strictly increasing, maximizing $L$ and maximizing $\log L$ give the \emph{same}
$\theta$, and the product becomes a sum:
\[
  \ell(\theta) = \log L(\theta) = \sum_{i=1}^n \log p(x_i,y_i\given\theta).
\]
This is the \emph{log-likelihood}, the central object of the chapter.

\begin{whybox}
Why does the maximizer not change when we take logs? If $\theta^\star$ maximizes
$L$, then $L(\theta^\star) \ge L(\theta)$ for all $\theta$. Apply the increasing
function $\log$ to both sides: $\log L(\theta^\star) \ge \log L(\theta)$. So
$\theta^\star$ also maximizes $\log L$. The logarithm relabels the heights of the
landscape but never moves a peak. We get numerical stability and additive
structure for free.
\end{whybox}

\begin{example}[The coin, made precise]
\label{ch34:coin-mle}
Model: $y_i \sim \Bern(p)$ i.i.d., so $p(y_i\given p) = p^{y_i}(1-p)^{1-y_i}$. With
$k = \sum y_i$ heads out of $n$,
\[
  \ell(p) = \sum_{i=1}^n \big[y_i\log p + (1-y_i)\log(1-p)\big]
          = k\log p + (n-k)\log(1-p).
\]
Differentiate and set to zero:
\[
  \ell'(p) = \frac{k}{p} - \frac{n-k}{1-p} = 0
  \;\Longrightarrow\; k(1-p) = (n-k)p
  \;\Longrightarrow\; \hat p = \frac{k}{n}.
\]
With $k=14$, $n=20$ we get $\hat p = 0.7$, precisely the gut answer. The second
derivative is negative, confirming a maximum. Your instinct \emph{was} calculus
on a log-likelihood all along.
\end{example}

Figure~\ref{ch34:fig-loglik} plots $\ell(p)$ for this data. Notice it is a single
smooth hill -- one peak, no traps. That concavity is what makes the coin (and,
later, logistic regression) so well-behaved.

\begin{figure}[H]
\centering
\begin{tikzpicture}
\begin{axis}[
  width=11cm, height=6.4cm,
  xlabel={$p$ (probability of heads)},
  ylabel={log-likelihood $\ell(p)$},
  domain=0.02:0.98, samples=120,
  xmin=0, xmax=1,
  grid=both,
  title={Log-likelihood for 14 heads in 20 flips}
]
\addplot[densityred, thick] {14*ln(x) + 6*ln(1-x)};
\addplot[dashed, gray] coordinates {(0.7,-30) (0.7,-12.2)};
\node[anchor=south] at (axis cs:0.7,-12.0) {$\hat p = 0.7$};
\end{axis}
\end{tikzpicture}
\caption{The log-likelihood is a single concave hill peaking exactly at the
sample fraction $\hat p = k/n = 0.7$. Maximum likelihood walks uphill to this
peak.}
\label{ch34:fig-loglik}
\end{figure}

\begin{historybox}
The likelihood as a tool for inference is largely the creation of Ronald A.
Fisher in the 1910s and 1920s. He insisted on a sharp distinction that still
trips people up: probability runs forward from a fixed parameter to random data;
\emph{likelihood} runs backward from fixed data to a varying parameter. They use
the same formula $p(\text{data}\given\theta)$ but read it along different axes.
Fisher's word ``likelihood'' was deliberate -- he did not want people to think a
parameter has a probability (that is the Bayesian move, coming later).
\end{historybox}

\section{MLE and MAP as training}
\label{ch34:mle-map}

\subsection{Maximum likelihood is minimizing surprise}

Define the maximum likelihood estimate
\[
  \hat\theta_{\text{MLE}} = \arg\max_\theta \ell(\theta)
                          = \arg\max_\theta \sum_{i=1}^n \log p(\text{datum}_i\given\theta).
\]
Flip the sign and you get a quantity machine learners love:
\[
  \hat\theta_{\text{MLE}} = \arg\min_\theta \underbrace{\Big[-\tfrac{1}{n}\sum_{i=1}^n \log p(\text{datum}_i\given\theta)\Big]}_{\text{average negative log-likelihood}}.
\]
The bracketed term is a \emph{loss function}. So ``maximize likelihood'' and
``minimize loss'' are the same sentence said two ways. This is the bridge between
the probability vocabulary and the optimization vocabulary of deep learning.

\begin{happeningbox}
What is the average negative log-likelihood really measuring? It is the
\emph{cross-entropy} between the empirical data distribution and your model (see
the information theory chapter). Minimizing it is the same as minimizing the KL
divergence $\KL(\hat p_{\text{data}} \,\|\, p_\theta)$: you are pulling your
model as close as possible, in the information-theoretic sense, to the histogram
of the data you saw. Every cross-entropy loss in every neural network is this
exact quantity. ``Train the network'' = ``do maximum likelihood.''
\end{happeningbox}

\begin{connectionbox}
The leap from ``minimize KL to the empirical distribution'' to ``minimize KL to
the true $p^\star$'' is justified by the law of large numbers: as $n\to\infty$,
the empirical average $\frac1n\sum -\log p_\theta(\text{datum}_i)$ converges to
its expectation $\E_{p^\star}[-\log p_\theta]$, which is minimized at the
$\theta$ closest to $p^\star$. So MLE is consistent precisely \emph{because} of
the LLN you proved earlier. The whole edifice rests on that one limit theorem.
\end{connectionbox}

\subsection{Overfitting and the case for a prior}

MLE has a famous weakness. Flip a coin three times, see three heads. MLE says
$\hat p = 3/3 = 1$: the coin \emph{never} lands tails. That is absurd, and worse,
it is dangerous -- the model now assigns probability zero to an outcome that can
certainly happen, and a single tail in the future makes the likelihood collapse.
With little data, maximum likelihood believes the data too literally. It
overfits.

The fix is to bring in prior knowledge. Before seeing any flips you already
believe a coin is roughly fair; you should need real evidence to be talked out of
that. Bayes' theorem is exactly the machinery for updating a belief with
evidence:
\[
  \underbrace{p(\theta\given\mathcal{D})}_{\text{posterior}}
  \;\propto\;
  \underbrace{p(\mathcal{D}\given\theta)}_{\text{likelihood}}\;
  \underbrace{p(\theta)}_{\text{prior}}.
\]

\subsection{MAP: maximum a posteriori}

If, instead of the value that maximizes the likelihood, we take the value that
maximizes the \emph{posterior}, we get the maximum a posteriori estimate:
\[
  \hat\theta_{\text{MAP}}
  = \arg\max_\theta \big[\log p(\mathcal{D}\given\theta) + \log p(\theta)\big]
  = \arg\max_\theta \big[\ell(\theta) + \log p(\theta)\big].
\]
The prior contributes an \emph{additive penalty} $\log p(\theta)$ to the
log-likelihood. This single line unifies two worlds.

\begin{example}[The coin, regularized]
\label{ch34:coin-map}
Put a $\Bet(\alpha,\beta)$ prior on $p$: $p(p)\propto p^{\alpha-1}(1-p)^{\beta-1}$.
The Beta is the conjugate prior for the Bernoulli (the posterior is again Beta),
which is why it is the natural choice. The posterior is
$\Bet(\alpha+k,\ \beta+n-k)$, and its mode (the MAP estimate) is
\[
  \hat p_{\text{MAP}} = \frac{k + \alpha - 1}{n + \alpha + \beta - 2}.
\]
With a mild fair-coin prior $\alpha=\beta=2$ and the three-heads disaster
$k=n=3$,
\[
  \hat p_{\text{MAP}} = \frac{3 + 1}{3 + 2} = 0.8,
\]
not the reckless $1.0$ of MLE. The prior acts like two phantom flips, one head and
one tail, added to the data. The Beta hyperparameters are ``pseudo-counts'': prior
evidence measured in the same currency as real evidence.
\end{example}

\begin{aestheticbox}
MAP reveals that \emph{regularization is a prior in disguise}. The famous
$L^2$ penalty (ridge regression, weight decay) $\lambda\norm{\theta}^2$ is
exactly $-\log p(\theta)$ for a Gaussian prior
$\theta\sim\Normal(0,\sigma^2 I)$ with $\lambda = 1/(2\sigma^2)$. The $L^1$
penalty (lasso) is a Laplace prior. So every time a practitioner ``adds
regularization to prevent overfitting,'' they are -- whether they know it or
not -- declaring a prior belief that the parameters are small. Weight decay is
Bayes.
\end{aestheticbox}

\begin{whybox}
Why does $L^2$ regularization equal a Gaussian prior? A Gaussian density is
$p(\theta)\propto \exp(-\norm{\theta}^2/(2\sigma^2))$, so
$\log p(\theta) = -\norm{\theta}^2/(2\sigma^2) + \text{const}$. The MAP objective
$\ell(\theta) + \log p(\theta)$ therefore equals
$\ell(\theta) - \tfrac{1}{2\sigma^2}\norm{\theta}^2 + \text{const}$. Maximizing
this is the same as minimizing $-\ell(\theta) + \lambda\norm{\theta}^2$ with
$\lambda=1/(2\sigma^2)$: ordinary penalized loss. The penalty strength and the
prior's tightness are two names for the same dial.
\end{whybox}

\begin{trapbox}
MAP gives a single point, the \emph{peak} of the posterior. The full Bayesian
answer is the entire posterior \emph{distribution}, which also carries the width
(your uncertainty). For a skewed posterior the mode can sit far from the mean and
be a poor summary. MAP is Bayes' theorem used as a regularizer, not Bayes' theorem
used to quantify uncertainty. Do not confuse the two.
\end{trapbox}

\section{Naive Bayes: a generative classifier}
\label{ch34:naive-bayes}

\subsection{Turning Bayes' theorem into a classifier}

Our first full-fledged learning algorithm. We want $p(y\given x)$ -- the
probability of class $y$ given features $x$. Bayes' theorem flips it:
\[
  p(y\given x) = \frac{p(x\given y)\,p(y)}{p(x)} \;\propto\; p(x\given y)\,p(y).
\]
To classify, pick the class with the largest right-hand side; the denominator
$p(x)$ is the same for every class and drops out. We need two ingredients: the
\emph{class prior} $p(y)$ (how common each class is) and the \emph{class-conditional}
$p(x\given y)$ (what features look like within each class).

The catch: $x$ is often high-dimensional -- thousands of word counts in an email --
and $p(x\given y)$ is a joint distribution over all of them, hopeless to estimate.

\subsection{The naive assumption}

Here is the bold, ``naive'' move: assume the features are
\emph{conditionally independent given the class}:
\[
  p(x\given y) = \prod_{j=1}^d p(x_j \given y).
\]
A monstrous joint distribution collapses into a product of $d$ one-dimensional
distributions, each trivial to estimate by counting. For spam: assume that, once
you know an email is spam, the appearance of ``viagra'' tells you nothing extra
about whether ``free'' also appears. This is false -- those words travel together --
yet the classifier works astonishingly well in practice.

Figure~\ref{ch34:fig-nbtree} shows the generative story as a tree: nature first
picks the class with probability $p(y)$, then -- treating the words as
independent given the class -- emits each word with its own class-conditional
probability.

\begin{figure}[H]
\centering
\begin{tikzpicture}[
  every node/.style={font=\small},
  level 1/.style={sibling distance=34mm, level distance=14mm},
  level 2/.style={sibling distance=17mm, level distance=14mm},
  edge from parent/.style={draw, -{Latex[length=2mm]}}
]
\node[draw, rounded corners, fill=blue!5] {class $y$}
  child {node[draw, rounded corners, fill=red!5] {spam}
    child {node {``free''?}}
    child {node {``viagra''?}}
    edge from parent node[left, pos=0.5] {$p(\texttt{spam})$}}
  child {node[draw, rounded corners, fill=green!5] {ham}
    child {node {``free''?}}
    child {node {``viagra''?}}
    edge from parent node[right, pos=0.5] {$p(\texttt{ham})$}};
\end{tikzpicture}
\caption{The generative tree behind naive Bayes. First a class is drawn from the
prior $p(y)$; then, \emph{given} the class, each word is emitted independently
with probability $p(x_j\given y)$. The independence is the ``naive'' part: the
branches for ``free'' and ``viagra'' do not talk to each other once the class is
fixed.}
\label{ch34:fig-nbtree}
\end{figure}

\begin{example}[Spam filtering by counting]
\label{ch34:spam}
Bag-of-words model: $x_j = 1$ if word $j$ appears. By MLE the parameters are pure
counts:
\[
  \hat p(y=\texttt{spam}) = \frac{\#\text{spam emails}}{\#\text{emails}},
  \qquad
  \hat p(x_j=1\given y) = \frac{\#\{y\text{-emails with word } j\}}{\#\,y\text{-emails}}.
\]
A new email with word-presence vector $x$ is scored by
\[
  \log\frac{p(\texttt{spam}\given x)}{p(\texttt{ham}\given x)}
  = \log\frac{p(\texttt{spam})}{p(\texttt{ham})}
   + \sum_{j} \log\frac{p(x_j\given\texttt{spam})}{p(x_j\given\texttt{ham})}.
\]
Each word casts an additive vote on the log-odds scale, and the votes simply add
up. Classify as spam when the total exceeds zero.
\end{example}

\begin{trapbox}
\textbf{The zero-frequency trap.} If a word never appeared in spam during
training, $\hat p(x_j=1\given\texttt{spam}) = 0$, and a single such word drives
the whole product -- and the spam probability -- to exactly zero, vetoing all
other evidence. The cure is \emph{Laplace smoothing}: add a pseudo-count to every
tally,
\[
  \hat p(x_j=1\given y) = \frac{\#\{\dots\} + 1}{\#\,y\text{-emails} + 2}.
\]
This is not an ad hoc patch: it is precisely the MAP estimate under a uniform
$\Bet(2,2)$ prior. The same phantom counts from Example~\ref{ch34:coin-map}
return to save us. Theory and the engineering hack are the same thing.
\end{trapbox}

\begin{paradoxbox}
\textbf{Naive Bayes is wrong yet wins.} Its independence assumption is almost
never true, so its probability \emph{estimates} are badly miscalibrated --
typically pushed toward $0$ or $1$ with false confidence. How can a model with
wrong probabilities classify well? Because classification needs only the
\emph{ranking} of classes to be correct, not the probabilities. As long as the
true class ends up with the largest score, the inflated magnitudes do not matter
for the decision. Naive Bayes can be a great classifier and a terrible
probability estimator at the same time -- a sharp reminder to separate ``which
class'' from ``how sure.''
\end{paradoxbox}

\section{Logistic regression as a probabilistic model}
\label{ch34:logistic}

\subsection{From a line to a probability}

Naive Bayes models how the features are generated. Logistic regression makes the
opposite bet: model $p(y\given x)$ \emph{directly}, without ever describing how
$x$ arises. Start with a linear score $z = w^\top x + b$, which ranges over all of
$\R$, and squash it into $(0,1)$ with the \emph{logistic} (sigmoid) function:
\[
  \sigma(z) = \frac{1}{1+e^{-z}}, \qquad
  p(y=1\given x) = \sigma(w^\top x + b).
\]
The sigmoid (Figure~\ref{ch34:fig-sigmoid}) is a smooth switch: very negative
scores map near $0$, very positive scores near $1$, and the steep middle is the
decision boundary.

\begin{figure}[H]
\centering
\begin{tikzpicture}
\begin{axis}[
  width=11cm, height=5.6cm,
  xlabel={score $z = w^\top x + b$},
  ylabel={$p(y=1\given x)$},
  domain=-7:7, samples=160,
  ymin=-0.05, ymax=1.05, xmin=-7, xmax=7,
  grid=both,
  title={The logistic (sigmoid) link}
]
\addplot[distblue, very thick] {1/(1+exp(-x))};
\addplot[dashed, gray] coordinates {(-7,0.5) (7,0.5)};
\addplot[dashed, gray] coordinates {(0,-0.05) (0,1.05)};
\end{axis}
\end{tikzpicture}
\caption{The sigmoid turns an unbounded linear score into a probability. At
$z=0$ the model is maximally uncertain ($p=0.5$); the slope there is steepest,
so small changes near the boundary swing the prediction most.}
\label{ch34:fig-sigmoid}
\end{figure}

\subsection{Why the sigmoid, of all functions?}

It is not arbitrary. The log-odds (logit) is the natural inverse:
\[
  \log\frac{p(y=1\given x)}{p(y=0\given x)} = w^\top x + b.
\]
Logistic regression is the assumption that the \emph{log-odds are linear in the
features}. And recall the naive-Bayes spam score: it too was linear in the
features on the log-odds scale. That is no coincidence.

\begin{connectionbox}
If the class-conditionals $p(x\given y)$ are Gaussians with a shared covariance,
Bayes' theorem produces \emph{exactly} a logistic form for $p(y\given x)$. So
Gaussian naive Bayes and logistic regression can describe the same boundary --
but they fit it differently. Naive Bayes estimates the Gaussians (generative);
logistic regression optimizes the boundary directly (discriminative). This is
the cleanest example of the generative/discriminative pair in
Section~\ref{ch34:gen-disc}.
\end{connectionbox}

\subsection{Training: maximum likelihood again}

With $y_i\in\{0,1\}$ and $\hat p_i = \sigma(w^\top x_i + b)$, each label is
Bernoulli, so the log-likelihood is
\[
  \ell(w,b) = \sum_{i=1}^n \big[y_i\log\hat p_i + (1-y_i)\log(1-\hat p_i)\big].
\]
Negate it and you have the \emph{binary cross-entropy} loss -- the workhorse loss
of binary classifiers everywhere. Maximizing this likelihood is training the
classifier; once more, learning is MLE.

There is no closed form for $\hat w$ this time, but the news is good: $\ell$ is
\emph{concave} in $(w,b)$, so gradient ascent cannot get stuck in a bad local
optimum. The gradient is strikingly clean:
\[
  \frac{\partial \ell}{\partial w} = \sum_{i=1}^n (y_i - \hat p_i)\,x_i.
\]
Each example pushes the weights in the direction of its features, scaled by the
\emph{residual} $y_i - \hat p_i$ -- how wrong the current prediction was.
Confident-and-correct examples ($\hat p_i\approx y_i$) contribute almost nothing;
the surprises do the teaching.

\begin{whybox}
Why is that gradient so simple? Two derivatives conspire. The sigmoid satisfies
$\sigma'(z) = \sigma(z)(1-\sigma(z))$, and the cross-entropy's derivative with
respect to $\hat p$ has a matching $\hat p(1-\hat p)$ in the denominator. They
cancel exactly, leaving the bare residual $y_i-\hat p_i$. This cancellation is
not luck -- it is the signature of \emph{matching the loss to the link} via the
exponential-family/canonical-link theory of generalized linear models. Use the
right loss for your link and the gradient always reduces to (truth minus
prediction) times input.
\end{whybox}

\begin{applicationbox}
Logistic regression is still, in the 2020s, one of the most deployed models on
Earth: credit scoring, click-through prediction, clinical risk scores, fraud
detection. Why, in the age of deep nets? It is fast, its coefficients are
interpretable (each $w_j$ is a log-odds change per unit feature), it produces
genuine probabilities, and it is the final layer of essentially every neural
classifier -- a deep net is logistic regression on top of learned features. Master
this and you have mastered the output end of modern deep learning.
\end{applicationbox}

\section{Generative vs.\ discriminative}
\label{ch34:gen-disc}

We have now met both species, so let us name the divide explicitly. A
\emph{generative} model learns the full joint $p(x,y) = p(x\given y)\,p(y)$ -- it
can, in principle, dream up new $(x,y)$ pairs. A \emph{discriminative} model
learns only the conditional $p(y\given x)$ -- it can answer questions about $y$
but cannot generate $x$.

\begin{center}
\begin{tabular}{@{}lll@{}}
\toprule
 & \textbf{Generative} & \textbf{Discriminative} \\
\midrule
Models & $p(x,y)$ & $p(y\given x)$ \\
Examples & naive Bayes, GMM, LDA & logistic regression, neural nets \\
Can sample new $x$? & yes & no \\
Small-data behavior & often better & needs more data \\
Large-data accuracy & can plateau & usually higher \\
\bottomrule
\end{tabular}
\end{center}

\begin{happeningbox}
There is a deep result behind this table (Ng and Jordan, 2002). A
generative--discriminative pair like Gaussian naive Bayes and logistic regression
converge to \emph{different} answers, and the generative model reaches its
(higher) asymptotic error faster -- it needs fewer examples -- while the
discriminative model reaches a lower asymptotic error but needs more data to get
there. The curves cross. Generative wins in the small-data regime; discriminative
wins once data is plentiful. This is not folklore; it is a theorem about the
price you pay for the generative model's stronger assumptions.
\end{happeningbox}

\section{Gaussian mixtures and the EM algorithm}
\label{ch34:gmm-em}

\subsection{When the labels are hidden}

Everything so far was \emph{supervised}: we saw labels $y_i$. Now suppose we see
only the features -- a cloud of points -- and suspect they come from a few distinct
groups (customer segments, cell types, signal sources). This is
\emph{clustering}, and the probabilistic way to think about it is a
\emph{mixture model}: each point was generated by first secretly choosing a group
$z\in\{1,\dots,K\}$, then drawing $x$ from that group's distribution.

For continuous data the natural choice is the \emph{Gaussian mixture}:
\[
  p(x) = \sum_{k=1}^K \pi_k\,\Normal(x \given \mu_k, \Sigma_k),
\]
a weighted blend of $K$ bell curves with mixing weights $\pi_k \ge 0$,
$\sum_k\pi_k = 1$. The hidden group label $z$ is a \emph{latent variable}: it was
part of the generative story but is missing from the data.

\begin{figure}[H]
\centering
\begin{tikzpicture}
\begin{axis}[
  width=11cm, height=5.6cm,
  xlabel={$x$}, ylabel={density},
  domain=-6:10, samples=180,
  ymin=0, grid=both, legend pos=north east,
  title={A two-component Gaussian mixture}
]
\addplot[dashed, sourceblue, thick] {0.4*exp(-(x-0)^2/2)/sqrt(2*pi)};
\addplot[dashed, quotientorange, thick] {0.6*exp(-(x-5)^2/(2*1.5^2))/(1.5*sqrt(2*pi))};
\addplot[densityred, very thick] {0.4*exp(-(x-0)^2/2)/sqrt(2*pi) + 0.6*exp(-(x-5)^2/(2*1.5^2))/(1.5*sqrt(2*pi))};
\legend{component 1, component 2, mixture}
\end{axis}
\end{tikzpicture}
\caption{The solid curve is the observed density; the dashed curves are the
hidden components. Clustering is the problem of recovering the dashed curves --
and which point came from which -- when you can see only the solid one.}
\label{ch34:fig-gmm}
\end{figure}

In two dimensions the same picture becomes a cloud of points
(Figure~\ref{ch34:fig-scatter}). Our eyes group it instantly into two blobs; the
EM algorithm below is a precise procedure for doing what the eye does, and for
attaching a probability to each membership rather than a hard verdict.

\begin{figure}[H]
\centering
\begin{tikzpicture}
\begin{axis}[
  width=8.4cm, height=6.0cm,
  xlabel={$x_1$}, ylabel={$x_2$},
  xmin=-4, xmax=8, ymin=-4, ymax=8,
  grid=both,
  title={Two clusters, labels hidden}
]
\addplot[only marks, mark=*, mark size=1.1pt, sourceblue] coordinates {
  (-0.4,0.6) (0.7,-0.3) (-1.1,0.2) (0.3,1.2) (1.0,0.4) (-0.6,-0.8)
  (0.1,-0.5) (-1.4,0.9) (0.9,1.1) (-0.2,-1.2) (1.3,-0.1) (-0.9,1.4)
  (0.5,0.0) (-0.3,-0.4) (0.8,-1.0) (-1.0,-0.2)};
\addplot[only marks, mark=*, mark size=1.1pt, quotientorange] coordinates {
  (4.6,4.4) (5.3,3.7) (3.9,5.2) (5.0,4.9) (4.2,3.6) (5.6,4.1)
  (4.8,5.5) (3.7,4.0) (5.4,5.0) (4.1,4.7) (5.1,3.9) (4.5,5.3)
  (3.8,4.5) (5.2,4.6) (4.9,3.8) (4.3,4.9)};
\end{axis}
\end{tikzpicture}
\caption{The same mixture in two dimensions: a cloud the eye splits into two
groups at a glance. Clustering automates that judgment and, with a Gaussian
mixture, reports \emph{how sure} it is about each point near the gap.}
\label{ch34:fig-scatter}
\end{figure}

\subsection{The chicken-and-egg problem and the EM idea}

Fitting a GMM by maximum likelihood faces a circular obstacle. \emph{If} we knew
each point's group, estimating each Gaussian would be trivial averaging. \emph{If}
we knew the Gaussians, assigning each point to its most probable group would be
easy. But we know neither. Expectation--Maximization (EM) breaks the circle by
guessing, then alternating:

\begin{itemize}
  \item \textbf{E-step.} Given the current Gaussians, compute for each point a
  \emph{soft} assignment -- the posterior probability it belongs to group $k$,
  called the \emph{responsibility}
  \[
    r_{ik} = \frac{\pi_k\,\Normal(x_i\given\mu_k,\Sigma_k)}
                  {\sum_{j=1}^K \pi_j\,\Normal(x_i\given\mu_j,\Sigma_j)}.
  \]
  This is just Bayes' theorem for the hidden label $z$.
  \item \textbf{M-step.} Given the responsibilities, re-estimate each Gaussian by
  \emph{weighted} maximum likelihood:
  \[
    \mu_k = \frac{\sum_i r_{ik} x_i}{\sum_i r_{ik}}, \quad
    \Sigma_k = \frac{\sum_i r_{ik}(x_i-\mu_k)(x_i-\mu_k)^\top}{\sum_i r_{ik}}, \quad
    \pi_k = \frac{1}{n}\sum_i r_{ik}.
  \]
\end{itemize}
Iterate E, M, E, M until the parameters stop moving.

\begin{intuitionbox}
EM is a guided dance between belief and evidence. The E-step says ``given what I
currently believe the clusters are, how should I split the credit for each
point?'' The M-step says ``given that split of credit, what clusters best explain
it?'' Each step can only \emph{increase} the likelihood, so the dance always
climbs -- though, unlike logistic regression, possibly to a local rather than the
global peak.
\end{intuitionbox}

\begin{whybox}
Why must each round increase the likelihood? EM does not maximize the awkward
log-likelihood $\log\sum_z p(x,z\given\theta)$ directly. Instead it maximizes a
\emph{lower bound} $\mathcal{Q}$ built using the current responsibilities (the
bound comes from Jensen's inequality applied to the concave $\log$). The E-step
makes the bound \emph{touch} the true log-likelihood at the current $\theta$; the
M-step pushes the bound -- and hence the touching log-likelihood -- uphill. Climbing
a bound that touches the true objective guarantees the true objective never
decreases. EM is hill-climbing where you cannot see the hill, only a tangent
surface you keep re-drawing.
\end{whybox}

\begin{connectionbox}
The familiar $k$-means algorithm is EM's stern cousin. Replace the soft
responsibilities with \emph{hard} assignments (each point to its single nearest
center) and fix all $\Sigma_k = \sigma^2 I$ with $\sigma\to 0$, and the EM updates
\emph{become} the $k$-means updates. So $k$-means is a Gaussian mixture with
spherical clusters and zero uncertainty about membership. Knowing this tells you
exactly when $k$-means fails: elongated or overlapping clusters, where the hard,
spherical assumption is plainly wrong.
\end{connectionbox}

\section{A first look at sampling and MCMC}
\label{ch34:mcmc}

\subsection{Why sampling is the engine room of Bayesian ML}

The honest Bayesian answer to any question is an expectation over the posterior:
the predictive mean is $\E_{\theta\sim p(\theta\given\mathcal{D})}[\,\cdot\,]$, a
credible interval is a posterior quantile, and so on. But the posterior
\[
  p(\theta\given\mathcal{D}) = \frac{p(\mathcal{D}\given\theta)\,p(\theta)}{p(\mathcal{D})},
  \qquad p(\mathcal{D}) = \int p(\mathcal{D}\given\theta)\,p(\theta)\,\dd\theta,
\]
hides a usually intractable integral in its denominator (the
\emph{evidence}). For all but a handful of conjugate models we cannot write the
posterior down. The escape: never compute the integral; instead draw
\emph{samples} $\theta^{(1)},\dots,\theta^{(m)}$ from the posterior and
approximate any expectation by an average,
\[
  \E[g(\theta)\given\mathcal{D}] \approx \frac{1}{m}\sum_{t=1}^m g(\theta^{(t)}),
\]
which is valid by the law of large numbers. The whole problem becomes: how do you
sample from a distribution you can only evaluate \emph{up to} the unknown constant
$p(\mathcal{D})$?

\subsection{Metropolis--Hastings: a random walk that obeys}

Markov chain Monte Carlo (MCMC) builds a Markov chain (Chapter on Markov chains)
whose stationary distribution \emph{is} the posterior. Run the chain long enough
and its states become samples from the target. The Metropolis--Hastings recipe:

\begin{enumerate}
  \item From the current $\theta$, \emph{propose} a nearby $\theta'$ (e.g.\
  $\theta' = \theta + \Normal(0,\tau^2)$).
  \item Compute the acceptance ratio
  $\displaystyle \alpha = \min\!\left(1,\ \frac{p(\mathcal{D}\given\theta')\,p(\theta')}
       {p(\mathcal{D}\given\theta)\,p(\theta)}\right).$
  \item With probability $\alpha$ move to $\theta'$; otherwise stay at $\theta$.
  Record the current state and repeat.
\end{enumerate}

\begin{happeningbox}
The magic is in step 2: the troublesome evidence $p(\mathcal{D})$ appears in both
numerator and denominator of the \emph{ratio} and cancels. We only ever need the
unnormalized posterior $p(\mathcal{D}\given\theta)p(\theta)$, which we can always
evaluate. The rule ``always climb to higher-probability proposals, and sometimes
climb down'' is what makes the chain's stationary distribution match the target
exactly -- this is the \emph{detailed balance} condition from the Markov chain
chapter. The occasional downhill moves are essential: a pure hill-climber would
collapse to the mode and never explore, never sample.
\end{happeningbox}

Figure~\ref{ch34:fig-mcmc} shows what such a chain looks like in practice: a
jagged walk that drifts from its arbitrary start into the bulk of the posterior
and then wanders there, each recorded state a (correlated) draw.

\begin{figure}[H]
\centering
\begin{tikzpicture}
\begin{axis}[
  width=12cm, height=5.0cm,
  xlabel={iteration $t$}, ylabel={state $\theta^{(t)}$},
  xmin=0, xmax=80, ymin=-1, ymax=6,
  grid=both,
  title={A Metropolis--Hastings trace (burn-in, then mixing)}
]
\addplot[distblue, thick] coordinates {
  (0,5.6)(1,5.6)(2,5.1)(3,4.7)(4,4.7)(5,4.2)(6,3.9)(7,3.9)(8,3.5)(9,3.1)
  (10,3.1)(11,2.8)(12,2.4)(13,2.6)(14,2.1)(15,2.1)(16,1.9)(17,2.3)(18,2.0)(19,1.7)
  (20,2.2)(21,2.4)(22,2.0)(23,1.6)(24,1.9)(25,2.5)(26,2.1)(27,1.8)(28,2.3)(29,2.0)
  (30,1.5)(31,1.9)(32,2.4)(33,2.7)(34,2.2)(35,1.8)(36,2.1)(37,2.6)(38,2.3)(39,1.7)
  (40,2.0)(41,2.5)(42,2.2)(43,1.6)(44,1.9)(45,2.4)(46,2.8)(47,2.3)(48,1.8)(49,2.1)
  (50,2.6)(51,2.2)(52,1.7)(53,2.0)(54,2.5)(55,2.1)(56,1.6)(57,2.3)(58,2.7)(59,2.2)
  (60,1.8)(61,2.1)(62,2.4)(63,2.9)(64,2.3)(65,1.9)(66,2.2)(67,2.6)(68,2.1)(69,1.7)
  (70,2.0)(71,2.5)(72,2.2)(73,1.8)(74,2.1)(75,2.4)(76,2.0)(77,1.6)(78,2.3)(79,2.1)};
\addplot[densityred, dashed, thick] coordinates {(15,-1)(15,6)};
\node[anchor=south, densityred] at (axis cs:15,5.2) {\small end of burn-in};
\end{axis}
\end{tikzpicture}
\caption{A one-dimensional Metropolis--Hastings trace. The chain starts high
(an arbitrary initial guess), descends through a transient \emph{burn-in}, then
oscillates around the posterior bulk near $\theta\approx 2$. Only the
post-burn-in states (right of the dashed line) are kept as samples, and even
those are correlated neighbors.}
\label{ch34:fig-mcmc}
\end{figure}

\begin{trapbox}
MCMC samples are \emph{correlated} (each state is a small step from the last) and
the early ``burn-in'' states reflect the arbitrary start, not the posterior. So
you discard an initial segment and remember that $m$ correlated draws carry less
information than $m$ independent ones. The proposal width $\tau$ is a delicate
dial: too small and the chain inches along, exploring nothing; too large and
almost every proposal is rejected, so the chain freezes in place. Diagnosing
whether a chain has ``mixed'' -- truly explored the posterior -- is a craft unto
itself. The Monte Carlo chapter develops these ideas in full.
\end{trapbox}

\begin{historybox}
The Metropolis algorithm was born in 1953 at Los Alamos, on the MANIAC computer,
in a paper by Metropolis, the Rosenbluths, and the Tellers -- aimed at simulating
the equilibrium of interacting particles, not statistics at all. W.K.\ Hastings
generalized it in 1970. For decades it lived in physics; only in the early 1990s,
when Gelfand, Smith, and others connected it to Bayesian inference, did it
detonate across statistics and machine learning. It is routinely listed among the
most influential algorithms of the twentieth century -- a physics trick that
quietly became the engine of modern Bayesian computation.
\end{historybox}

\section{Calibration and uncertainty}
\label{ch34:calibration}

\subsection{A confident model can be a liar}

A model that outputs probabilities makes an implicit promise. When it says ``$70\%$
chance of rain,'' it is promising that across all the days it makes that call, it
rains about $70\%$ of the time. A model that keeps this promise is
\emph{calibrated}. Calibration is a property of the probabilities, entirely
separate from accuracy: a model can be accurate but overconfident (think of
naive Bayes pinning probabilities to $0$ or $1$), or modest but perfectly honest.

\begin{focusbox}
\textbf{Accuracy} asks: is the top-scoring class usually right? \textbf{Calibration}
asks: do the probabilities mean what they say? A self-driving car or a cancer
screen needs both -- a wrong-and-confident model is far more dangerous than a
wrong-and-uncertain one, because the confidence is what downstream decisions
trust.
\end{focusbox}

\subsection{Reading a reliability diagram}

To check calibration, bin predictions by their stated probability and, in each
bin, plot the \emph{actual} fraction of positives against the \emph{average
predicted} probability. Perfect calibration lies on the $45^\circ$ line. A curve
sagging below it means the model is overconfident; bulging above means
underconfident (Figure~\ref{ch34:fig-calib}).

\begin{figure}[H]
\centering
\begin{tikzpicture}
\begin{axis}[
  width=9.6cm, height=6.4cm,
  xlabel={mean predicted probability},
  ylabel={observed fraction positive},
  xmin=0, xmax=1, ymin=0, ymax=1,
  grid=both, legend pos=north west,
  title={Reliability diagram}
]
\addplot[gray, dashed, thick] coordinates {(0,0) (1,1)};
\addlegendentry{perfectly calibrated}
\addplot[densityred, very thick, mark=*] coordinates {
  (0.1,0.04) (0.3,0.16) (0.5,0.34) (0.7,0.55) (0.9,0.78)};
\addlegendentry{overconfident model}
\end{axis}
\end{tikzpicture}
\caption{The overconfident model's curve sits below the diagonal: when it claims
$0.9$, positives occur only $\approx 0.78$ of the time. Such a curve is the
fingerprint of a model that needs recalibrating.}
\label{ch34:fig-calib}
\end{figure}

\begin{applicationbox}
Modern deep neural networks are notoriously overconfident -- far more so than the
simpler logistic regression they end in (Guo et al., 2017). Practical fixes are
cheap. \emph{Temperature scaling} divides the pre-sigmoid scores by a single
learned constant $T>1$, softening the probabilities without changing the
predicted class or the accuracy at all -- it only adjusts the confidence. \emph{Platt
scaling} (a tiny logistic regression on the scores) and \emph{isotonic regression}
do the same job nonparametrically. In any setting where the probability itself
drives a decision -- medicine, finance, autonomy -- calibration is not optional
polish; it is a safety requirement.
\end{applicationbox}

\subsection{Two kinds of uncertainty}

Finally, a distinction that matters whenever a model must say ``I don't know.''
\emph{Aleatoric} uncertainty is irreducible noise in the world -- a fair coin is
genuinely $50/50$ no matter how much data you gather. \emph{Epistemic}
uncertainty is your ignorance about the model itself, and \emph{shrinks} with more
data. The Bayesian posterior tracks epistemic uncertainty automatically: a
parameter seen by little data has a wide posterior, and that width flows through
into wide predictive intervals -- the model knows what it does not know.

\begin{aestheticbox}
This is the deepest reason to think probabilistically about machine learning.
A point-estimate model (plain MLE, an ordinary neural net) gives an answer and no
honest sense of how much to trust it. A probabilistic model gives a
\emph{distribution} -- it can report ``I am confident'' or ``I am guessing,'' and
out-of-distribution inputs, where it should be least sure, are exactly where a
calibrated probabilistic model widens its intervals. In high-stakes deployment,
knowing \emph{when not to trust the model} is worth as much as the prediction
itself. That self-knowledge is what probability buys you.
\end{aestheticbox}

\section{The thread, pulled tight}

Step back and look at the single thread running through the chapter. We posit a
probabilistic model $p(\text{data}\given\theta)$. We score parameters by the
likelihood, optionally tempered by a prior. Maximizing that score -- by calculus
when we are lucky, by gradient ascent or EM when we are not -- is what everyone
calls ``training.'' Bayes' theorem turns a generative model into a classifier
(naive Bayes) and turns priors into the regularizers practitioners reach for by
instinct (MAP). When integrals defeat us we sample our way around them (MCMC).
And throughout we insist the probabilities mean what they claim (calibration) and
report honest uncertainty. Maximum likelihood, Bayes, the normal distribution,
KL divergence, the law of large numbers, Markov chains -- every tool you built
earlier in this book reappears here, doing the real work of modern machine
learning. The methods change with the decade; this probabilistic skeleton does
not.

\section*{Exercises}
\addcontentsline{toc}{section}{Exercises}

\begin{exercise}[Poisson MLE]
You record the number of customer arrivals per hour over $n$ hours,
$y_1,\dots,y_n$, modeled as i.i.d.\ $\Pois(\lambda)$. Write the log-likelihood
$\ell(\lambda)$, differentiate, and show $\hat\lambda_{\text{MLE}} = \bar y$, the
sample mean. Verify it is a maximum via the second derivative.
\end{exercise}

\begin{exercise}[Regularization is a prior]
Show that minimizing the ridge objective
$-\ell(\theta) + \lambda\norm{\theta}^2$ is exactly MAP estimation under a
Gaussian prior $\theta\sim\Normal(0,\sigma^2 I)$. Express $\lambda$ in terms of
$\sigma^2$. Then state, in one sentence, what prior belief the lasso ($L^1$
penalty) encodes and why it tends to drive coefficients exactly to zero.
\end{exercise}

\begin{exercise}[Naive Bayes by hand]
A spam corpus has $60$ spam and $40$ ham emails. The word ``free'' appears in
$45$ of the spam and $8$ of the ham. Using Laplace smoothing (add one to each
count, two to each denominator), compute $\hat p(\text{spam})$,
$\hat p(\text{free}\given\text{spam})$, and
$\hat p(\text{free}\given\text{ham})$. For an email containing ``free'' (and
treating it as the only feature), compute the log-odds in favor of spam and give
the classification.
\end{exercise}

\begin{exercise}[Sigmoid and log-odds]
Starting from $p(y=1\given x) = \sigma(w^\top x + b)$ with
$\sigma(z) = 1/(1+e^{-z})$, prove algebraically that
$\log\frac{p(y=1\given x)}{p(y=0\given x)} = w^\top x + b$. Then show the identity
$\sigma'(z) = \sigma(z)(1-\sigma(z))$ and explain in one sentence how it makes the
logistic-regression gradient collapse to $(y_i - \hat p_i)x_i$.
\end{exercise}

\begin{exercise}[Cross-entropy is negative log-likelihood]
For binary labels $y_i\in\{0,1\}$ with predictions $\hat p_i$, write the Bernoulli
log-likelihood and show that its negative is the binary cross-entropy
$-\sum_i[y_i\log\hat p_i + (1-y_i)\log(1-\hat p_i)]$. Explain in one sentence
the sense in which minimizing it minimizes a KL divergence to the empirical
distribution.
\end{exercise}

\begin{exercise}[One EM update]
You fit a two-component GMM in one dimension. The current parameters are
$\pi_1=\pi_2=0.5$, $\mu_1=0$, $\mu_2=4$, and $\sigma_1=\sigma_2=1$. For the data
point $x=1$, compute the responsibilities $r_{11}$ and $r_{12}$ (the E-step).
Which component currently claims more of the point, and does that match your
intuition from the distance to each mean?
\end{exercise}

\begin{exercise}[Why the evidence cancels]
In Metropolis--Hastings the acceptance ratio uses
$p(\mathcal{D}\given\theta')p(\theta')$ over
$p(\mathcal{D}\given\theta)p(\theta)$. Explain precisely why we never need the
normalizing constant
$p(\mathcal{D}) = \int p(\mathcal{D}\given\theta)p(\theta)\,\dd\theta$,
and why this is what makes MCMC practical for models where that integral is
impossible.
\end{exercise}

\begin{exercise}[Calibration vs.\ accuracy, conceptual]
Construct a small concrete example of a binary classifier that is $90\%$ accurate
yet badly miscalibrated, and a second classifier that is only $60\%$ accurate yet
perfectly calibrated. (Describe the predicted probabilities and outcomes; numbers,
not prose.) Use your examples to explain why accuracy alone is the wrong metric
when a downstream decision depends on the probability itself.
\end{exercise}

\begin{exercise}[Generative vs.\ discriminative, proof-flavored]
Assume two classes with class-conditional densities
$p(x\given y=c) = \Normal(x\given\mu_c,\Sigma)$ sharing one covariance $\Sigma$,
and class prior $p(y=1)=\pi$. Using Bayes' theorem, show that $p(y=1\given x)$ has
the logistic form $\sigma(w^\top x + b)$, and give $w$ and $b$ in terms of
$\mu_0,\mu_1,\Sigma,\pi$. Comment on what this reveals about the relationship
between Gaussian naive Bayes and logistic regression.
\end{exercise}

\begin{exercise}[Aleatoric vs.\ epistemic]
Classify each source of uncertainty as primarily aleatoric or epistemic, with a
one-line justification: (a) the outcome of a future fair die roll; (b) a drug's
effect size estimated from only $12$ patients; (c) sensor readout noise on a
thermometer; (d) a self-driving model's prediction on a kind of road sign absent
from its training set. For the epistemic ones, state what would reduce the
uncertainty.
\end{exercise}
